
% a description of how the papers fit together and their relationship to the overall body of work in your topic, can be multiple paragraphs

%-------------------------------------------------------------------------
The selected literature collectively forms a cohesive narrative arc that mirrors the rapid maturation of AI agent security research. Specifically, these papers progress logically from theoretical threat modeling to empirical benchmarking, and ultimately to active defense mechanisms. While traditional LLM security research primarily focused on user-prompt injection, this body of work addresses a critical paradigm shift: vulnerabilities arising from the orchestration of external tool environments. By mapping this transition, the papers establish a new subfield focused on the inherent risks of delegating autonomous execution to loosely coupled, third-party servers, identifying the theoretical boundaries of context provider exploitation.

This theoretical framework is subsequently operationalized by the three benchmarking studies: MCP-SafetyBench \cite{zhang2025msb}, MCPSecBench \cite{yang2025mcpsecbench}, and MSB \cite{zhang2025msb}. While MCP-SafetyBench focuses heavily on host-side deception and deceptive server metadata, MCPSecBench expands the scope to evaluate data exfiltration and privilege escalation across diverse threat domains. Complementing both, MSB introduces the granular scoring methodologies required to quantify these failures precisely. Together, these three benchmarks provide a unified empirical consensus: state-of-the-art models inherently over-trust external servers, making them highly susceptible to tool poisoning and malicious context injection.

Finally, the introduction of the Model Contextual Integrity Protocol (MCIP)\cite{jing2025mcip} represents the necessary evolution from identifying flaws to mitigating them. MCIP directly addresses the systemic failures exposed by the aforementioned benchmarks by enforcing strict contextual boundaries between the reasoning host and external servers. In the broader landscape of multi-agent orchestration frameworks, such as CrewAI \cite{venkadesh2024crewai}, and Microsoft AutoGen\cite{wu2024autogen}, this collective body of work underscores that secure agentic deployment cannot rely solely on the underlying LLM's alignment. Instead, robust security requires architectural interventions at the protocol level to verify and isolate external context.